\emph{(i)} Each window $\mathfrak{A}_{F}$ is finite-dimensional and $\tau$-invariant, and on it $\tau_{t}$ acts through the phases $e^{-I_4\omega_{j}t}$, $j \in F$ (numbers read in $\mathrm{span}\{1, I_4\}$), so the orbit of any $x \in \mathfrak{A}_{F}$ is the continuous image of a subset of a compact torus and is relatively compact (Theorem \claimref{synthesis.point-spectrum-boost-admits-no-half-sided-translations}(i)). The union of the windows is dense in $\mathfrak{A}$ (Definition \claimref{obj.hawking.boundary-mode-algebra-windows-boost-dynamics}: $\mathfrak{A}$ is their C$^{*}$-inductive limit). Given $x \in \mathfrak{A}$ and $\varepsilon > 0$, pick $x'$ in a window with $\|x - x'\| < \varepsilon/3$; since every $\tau_{t}$ is isometric, $\|\tau_{t}(x) - \tau_{t}(x')\| < \varepsilon/3$ for all $t$, and finitely many balls of radius $\varepsilon/3$ cover the orbit of $x'$, so finitely many balls of radius $\varepsilon$ cover the orbit of $x$. The orbit is totally bounded, hence relatively compact in the Banach space $\mathfrak{A}$. For recurrence, fix $S > 0$ and consider $\tau_{nS}(x)$, $n = 1, 2, \dots$; they lie in a compact set, so two of them, $n < m$, satisfy $\|\tau_{mS}(x) - \tau_{nS}(x)\| < \varepsilon$, and applying the isometry $\tau_{-nS}$ gives $\|\tau_{(m - n)S}(x) - x\| < \varepsilon$ with $T = (m - n)S \ge S$. The points $\tau_{-nS}(x)$ give $T' \le -S$ in the same way.

\emph{(ii)} Suppose $\tau_{t}(\mathcal{B}) \subseteq \mathcal{B}$ for $t \ge 0$. Let $x \in \mathcal{B}$, $s > 0$ and $y = \tau_{-s}(x)$. For $\varepsilon > 0$, (i) gives $T \ge s$ with $\|\tau_{T}(y) - y\| < \varepsilon$, and $\tau_{T}(y) = \tau_{T - s}(x) \in \mathcal{B}$ because $T - s \ge 0$. So $y$ lies within $\varepsilon$ of $\mathcal{B}$ for every $\varepsilon$, and $y \in \mathcal{B}$ since $\mathcal{B}$ is closed. Thus $\tau_{-s}(\mathcal{B}) \subseteq \mathcal{B}$ for all $s > 0$, so $\tau_{t}(\mathcal{B}) \subseteq \mathcal{B}$ for every real $t$, and applying this to $t$ and $-t$ gives $\tau_{t}(\mathcal{B}) = \mathcal{B}$. The case $t \le 0$ is the same with $T' \le -s$. \emph{The GNS version.} Let $\omega$ be $\tau$-invariant with GNS triple $(\mathcal{H}, \pi, \Omega)$ and $V_{t}\pi(x)\Omega = \pi(\tau_{t}(x))\Omega$. Each $\pi(\mathfrak{A}_{F})\Omega$ is a finite-dimensional $V$-invariant subspace, and their union is dense, so $\mathcal{H}$ is spanned by eigenvectors of $V$: $V_{t} = \sum_{\mu}e^{-I_4\mu t}P_{\mu}$ with $\sum_{\mu}P_{\mu} = 1$, pure point spectrum. On unitaries of this form the strong topology is the product topology on the phases $(e^{-I_4\mu t})_{\mu}$, so the strong closure $K$ of $\{V_{t}\}$ is a closed subgroup of a product of circles, a compact abelian group. The strong closure of $\{V_{t} : t \ge 0\}$ is a closed subsemigroup of $K$, and a closed subsemigroup of a compact group is a subgroup; so for $s > 0$ there is a net $t_{\alpha} \ge 0$ with $V_{t_{\alpha}} \to V_{-s}$ strongly, hence strongly with adjoints (inversion is continuous on the unitary group in the strong topology). For $x \in \mathcal{B}$, $V_{t_{\alpha}}xV_{t_{\alpha}}^{*} \in \mathcal{B}$ is a bounded net converging strongly to $V_{-s}xV_{-s}^{*}$; on bounded sets the weak and $\sigma$-weak topologies agree, and a $\sigma$-weakly closed subspace meets every ball in a weakly closed set (Krein-Smulian; this and the subsemigroup fact are standard, cited from memory), so the limit lies in $\mathcal{B}$. The rest is as above.

\emph{(iii)} Under (G) and (P) the modular group is $\sigma_{s} = \tau_{\gamma s}$ for some real $\gamma \ne 0$ (clause (i) supplies the algebra and its boost dynamics $\tau$, (P) says that $\tau$ is how the local boost acts on the patch's modes, clause (iii) identifies the modular flow with that boost; the rescaling of the parameter, fixed in clause (ii) by the temperature, does not matter here; in the GNS form $\Delta^{I_4 s} = V_{\gamma s}$ and the GNS version of (ii) is used), and $\sigma_{s}(\mathfrak{A}_{c}) = \mathfrak{A}_{e^{\mp 2\pi s}c}$. Take the sign so that this is $\mathfrak{A}_{e^{-2\pi s}c}$ (the other is the same argument with $s$ replaced by $-s$). For $s \le 0$, $e^{-2\pi s}c \ge c$, so $\sigma_{s}(\mathfrak{A}_{c}) = \mathfrak{A}_{e^{-2\pi s}c} \subseteq \mathfrak{A}_{c}$ by the nesting. So $\mathcal{B} = \mathfrak{A}_{c}$, norm-closed as a C$^{*}$-subalgebra, is mapped into itself by $\tau_{t}$ for all $t$ of one sign, and (ii) gives $\sigma_{s}(\mathfrak{A}_{c}) = \mathfrak{A}_{c}$ for all $s$, that is $\mathfrak{A}_{e^{-2\pi s}c} = \mathfrak{A}_{c}$ for all $s$; as $s$ ranges over $\R$, $e^{-2\pi s}c$ ranges over all $c' > 0$. For nested diamonds $D_{c} \subseteq D$ sharing the future tip of $D$, the conformal Killing flow of $D$ moves the family monotonically (the conformal map from the wedge to the diamond carries the half null planes $\mathcal{N}_{c}$ to them, as in Hislop and Longo's treatment, cited from memory), so the same argument applies with $\mathfrak{A}_{c}$ the algebra of $D_{c}$. Only the phase action of $\tau$ is used: no property of the state enters.

\emph{The reading.} At a Killing horizon the store's modes sit on the fixed locus of the boost (Definition \claimref{obj.hawking.boundary-mode-algebra-windows-boost-dynamics}: one ladder pair per point of the blown-up singular boundary, with the boost acting by phases), so the flow acts on the algebra as a whole and moves no sub-patch; clause (iii) read that way is consistent with clause (i), and for a faithful state the modular group being $\tau_{-\beta s}$ is the $(\tau, \beta)$-KMS condition of clause (ii) (on windows this equivalence is part (ii) of Theorem \claimref{synthesis.local-boost-temperature-is-local-regularity}). The contrast in the check is the bilateral shift $S e_{k} = e_{k + 1}$ on $\ell^{2}(\mathbb{Z})$, unitary with Lebesgue spectrum, the model to which the dilations of the half line are unitarily equivalent: $\|S^{n}e_{0} - e_{0}\|^{2} = 2$ for every $n \ge 1$, and $S$ maps the closed span of $\{e_{k}\}_{k \ge 0}$ onto that of $\{e_{k}\}_{k \ge 1}$, a proper subspace, exactly the behaviour (ii) excludes for phase flows. The check illustrates and does not test (ii) or (iii): it verifies, on a two-mode window in the Jordan-Wigner representation, the canonical anticommutation relations and $\tau(a_{j}) = \bar z^{\omega_{j}}a_{j}$ for the step $z = (3 + 4I_4)/5$ with weights $(1, 2)$; $z = (2 + I_4)/(2 - I_4)$ is a quotient of non-associate Gaussian primes, so $z^{n} = 1$ would give $(2 + I_4)^{n} = (2 - I_4)^{n}$, impossible by unique factorisation in $\mathbb{Z}[I_4] = \mathrm{span}_{\mathbb{Z}}\{1, I_4\}$, and the flow is not periodic; still $|z^{61} - 1|^{2} < 10^{-3}$, the recurrence of (i) seen in exact arithmetic.
